% This source is also embedded in main.tex; compile independently with XeLaTeX/BibTeX.
\ifdefined\PaperMainDocument
  \newcommand{\EndSupplementDocument}{}
\else
% !TeX program = xelatex
% Chinese research manuscript using the Elsevier article class.
\documentclass[preprint,12pt]{elsarticle}
\usepackage[a4paper,margin=25mm]{geometry}
\usepackage[UTF8,fontset=none,scheme=plain]{ctex}
\usepackage{fontspec}
\IfFontExistsTF{Times New Roman}{\setmainfont{Times New Roman}}{\setmainfont{texgyretermes-regular.otf}[BoldFont=texgyretermes-bold.otf,ItalicFont=texgyretermes-italic.otf,BoldItalicFont=texgyretermes-bolditalic.otf]}
\IfFontExistsTF{SimSun}{\setCJKmainfont[AutoFakeBold=2,ItalicFont=SimSun,BoldItalicFont=SimSun]{SimSun}\setCJKmonofont{SimSun}}{\setCJKmainfont[BoldFont=FandolSong-Bold.otf,ItalicFont=FandolSong-Regular.otf,BoldItalicFont=FandolSong-Bold.otf]{FandolSong-Regular.otf}\setCJKmonofont{FandolSong-Regular.otf}}
\usepackage{amsmath,amssymb,unicode-math}
\setmathfont{latinmodern-math.otf}
\usepackage{etoolbox}
\usepackage{graphicx,booktabs,tabularx,array,caption,setspace,microtype,xurl,placeins,longtable,amsthm}
\biboptions{numbers,sort&compress}
\usepackage[hidelinks,unicode]{hyperref}
\hypersetup{pdftitle={补充材料：音乐游戏活动终分预测},pdfauthor={Tano},pdfsubject={Entertainment Computing 格式中文研究稿}}
\urlstyle{same}
\newcolumntype{L}{>{\raggedright\arraybackslash}X}
\newcolumntype{C}[1]{>{\centering\arraybackslash}p{#1}}
\captionsetup{font=small,labelfont=bf,justification=justified,singlelinecheck=false}
\renewcommand{\figurename}{图}
\renewcommand{\tablename}{表}
\renewcommand{\refname}{参考文献}
\abstracttitle{摘要}
\keywordtitle{\upshape 关键词}
\setstretch{1.15}
\setlength{\parindent}{2em}
\setlength{\parskip}{0.15em}
\setlength{\emergencystretch}{3em}
\hyphenpenalty=10000
\exhyphenpenalty=10000
\setlength{\textfloatsep}{12pt plus 2pt minus 2pt}
\setlength{\intextsep}{12pt plus 2pt minus 2pt}
\renewcommand{\topfraction}{0.92}
\renewcommand{\bottomfraction}{0.9}
\renewcommand{\textfraction}{0.05}
\renewcommand{\floatpagefraction}{0.8}
\setcounter{topnumber}{3}
\setcounter{totalnumber}{4}
\journal{Entertainment Computing}
% Show the manuscript page number without implying submission or acceptance.
\makeatletter
\patchcmd{\pprintMaketitle}{\vskip36pt}{\vskip12pt}{}{\errmessage{Title spacing patch failed}}
\def\ps@pprintTitle{\let\@oddhead\@empty\let\@evenhead\@empty
  \def\@oddfoot{\hfil\thepage\hfil}\let\@evenfoot\@oddfoot}
\makeatother

\newtheorem{proposition}{命题}
\renewcommand{\proofname}{证明}
\newcommand{\upfigure}[4]{\begin{figure}[!htbp]\centering\includegraphics[width=\linewidth,height=0.70\textheight,keepaspectratio]{figures/#1.pdf}\caption{#3}\label{#2}\end{figure}}
% Generated from finalized independent evidence; run scripts/bind_upgrade_evidence.py.
\newcommand{\BaseTest}{7.082}
\newcommand{\BaseVal}{6.672}
\newcommand{\EwmaTest}{6.671}
\newcommand{\EwmaVal}{6.502}
\newcommand{\PairTest}{6.540}
\newcommand{\PairVal}{6.556}
\newcommand{\DifferenceTest}{6.511}
\newcommand{\DlmTest}{6.672}
\newcommand{\PairDelta}{-0.130}
\newcommand{\PairAnchorIncrement}{0.030}
\newcommand{\DlmDelta}{0.001}
\newcommand{\PairTwoTest}{6.5397}
\newcommand{\PairLooLow}{-0.163}
\newcommand{\PairLooHigh}{-0.097}
\newcommand{\PairTopThree}{50.5}
\newcommand{\PairImproved}{11}
\newcommand{\TreeSharedFull}{6.912}
\newcommand{\TreeIndependentFull}{6.312}
\newcommand{\NnSharedFull}{10.949}
\newcommand{\NnHeadsFull}{10.198}
\newcommand{\NnIndependentFull}{9.314}
\newcommand{\CoverageEighty}{93.01}
\newcommand{\WidthEighty}{35.07}
\newcommand{\CoverageNinety}{96.90}
\newcommand{\WidthNinety}{47.64}
\newcommand{\JointEighty}{79.86}
\newcommand{\JointNinety}{90.47}
\newcommand{\ProspectiveRegistered}{2026-11-18 15:02:49.992 UTC}
\newcommand{\ProspectiveTarget}{20}
\newcommand{\ProspectiveEventCount}{0}

% Generated only from sealed evidence.
\newcommand{\RTIndependent}{5.918}
\newcommand{\RSIndependent}{4.977}
\newcommand{\RTQuarter}{5.954}
\newcommand{\RTPool}{6.816}
\newcommand{\RTRisk}{6.017}
\newcommand{\RTGated}{5.918}
\newcommand{\RTPoolDelta}{0.899}
\newcommand{\RTPoolLow}{-0.402}
\newcommand{\RTPoolHigh}{2.012}
\newcommand{\RTRiskDelta}{0.100}
\newcommand{\RTRiskLow}{-0.005}
\newcommand{\RTRiskHigh}{0.225}
\newcommand{\DContemp}{0.4349}
\newcommand{\DLagOne}{0.1549}
\newcommand{\DLagTwo}{-0.0174}
\newcommand{\DQMin}{-0.0010445}
\newcommand{\DCorrLow}{0.1431}
\newcommand{\DCorrHigh}{0.6752}
\newcommand{\DRatio}{7.48}
\newcommand{\DTransfer}{6.953329}
\newcommand{\DTransferLog}{0.008115}
\newcommand{\DTransferI}{5.340927}
\newcommand{\EWidthEighty}{35.20}
\newcommand{\EWinklerEighty}{958.57}
\newcommand{\EWidthNinety}{47.78}
\newcommand{\EWinklerNinety}{1239.23}
\newcommand{\ConfirmTarget}{80}
\newcommand{\ConfirmFloor}{60}
\newcommand{\ConfirmGlobalN}{4}
\newcommand{\ConfirmPairN}{72}
\newcommand{\ConfirmFuture}{0}


\begin{document}
\begin{frontmatter}
\title{补充材料：音乐游戏活动终分预测}
\author{Tano}
\end{frontmatter}
\appendix
\newcommand{\EndSupplementDocument}{\end{document}}
\fi
\section{补充材料：证据来源、完整菜单与审计范围}\label{sec:supplement}
本材料同时提供正文附录与独立版本，完整保留既有实验并补充训练期诊断。数值表使用原生LaTeX文本。只读绑定器核对封存哈希，并在\path{qa/upgrade_number_bindings.json}与\path{qa/round2_number_bindings.json}记录来源、字段、行条件、单位及舍入。图表的输入、代码与导出指纹由独立manifest记录。主文围绕研究问题解释，完整候选、预算及失败记录在此展开。

\subsection{证据层级与运行索引}
\begin{table}[!htbp]\centering
\caption{研究工件与结论范围}\label{tab:evidencelevels}
\small\begin{tabularx}{\linewidth}{p{3.0cm} L L}\toprule
工件 & 内容 & 解释范围\\\midrule
历史329行汇总 & 2026年11月15日首次菜单，14个协议与运行组合 & 同一方法可含种子和集成行，后续新方法另计\\
\path{fixed_replay_v2} & 65期查询、54种固定规则，19线基础预测逐例对齐 & 已探索测试上的完整机制与敏感性比较\\
\path{two_rank_learning_curve_v2} & 16/32/48/67期训练，12次树、48次神经拟合，56评分单元 & 同一14期488案例验证，局部自身特征和新基础配方\\
\path{matched_train_v3} & 34期18线共同支持相关及重采样 & 训练期描述与相关差诊断\\
\path{synthetic} & 16格、3种子、100序列、4种历史、9方法 & 指定生成模型中的机制与Monte Carlo误差\\
二轮B/C1 & 35期，15/10/10；12候选及18消融 & 训练期回溯外推，共用已发基础预测\\
二轮C2 & 3终点、3历史组、3种子；108拟合 & 12期查询全为35期子集，总参数和近似MAC分别控制\\
二轮D & 前25期校准、24格模拟、末10期事后迁移 & 受限机制诊断，整体运行末封存的迁移身份单列\\
二轮E & 独立重算、76+76条件、功效模型 & 固定损失与探索面板，0个未来活动\\
前瞻登记 & 原20期pilot及新80期独立确认计划 & 新计划待激活，未来精度留空\\\bottomrule
\end{tabularx}\end{table}

固定回放运行标识为\path{dcc493f82e83933d6963e2d39a02e6c624057ad6c0932ad24801efbd252b1f09}；其盲预测根清单SHA-256为\path{6c3a034e9e409c12230cb029f17a3268fec1f051fcce1a1eb34f1ee46390128a}。新拟合学习曲线运行标识为\path{ed9ea52183f47ef4ea566d41d5a52f4ff4661df5385957dae4865a29ae22ba67}，盲预测根清单为\path{45b87d8b9b9123eafde8fea0682b49620d6865e24c9a642032935d17ff1ee195}。评分代码、依赖及输出清单有独立分析标识。根目录保存的v1中断运行和较早学习曲线均独立归档，本文只绑定完成且核验的v2运行。

329行的统计键为\((\text{协议},\text{运行},\text{方法/种子/集成})\)，对应\path{data/supplement_all_329_results.csv}。其中14个协议与运行组合的计数沿用首次菜单；它描述该次归档范围。\path{data/supplement_historical_selections.csv}保留原选中方法和轮次，正文既有神经结果来自\path{data/main_models.csv}。本轮固定菜单的选择身份为\texttt{ewma\_reference}，后续评分没有新增测试最小值选择器。

\subsection{冻结数据指纹与案例成员关系}
\begingroup\footnotesize\setlength{\tabcolsep}{3pt}
\begin{longtable}{p{2.2cm} p{12.8cm}}
\caption{冻结数据完整SHA-256标识}\label{tab:datasetids}\\
\toprule
协议 & 数据指纹 \\
\midrule\endfirsthead
\toprule
协议 & 数据指纹 \\
\midrule\endhead
V1两线 & \path{80e354c3a5340df6724a40ef1df490f92c08007311d55849342de509eaf5ee50} \\
V2两线 & \path{d057beca6597dd28b59fde7da08f665ae74af41d96f940da04268c881864a215} \\
V2 19线 & \path{8ffd781fc9f4ebcad9ea691839777458e5e4c7cbeade58978dca15bce2440ad8} \\
V3两线 & \path{85f5924e0b90deb0e1a30f3aa27d2237640cca4d968d96612dda8fee1d9c72ec} \\
V3 19线 & \path{db4e9b9c84bd64a6de5ee30a090cfac83efbe85c2f6dd899304516e7d031cdf1} \\
\bottomrule
\end{longtable}
\endgroup

所有数据归属按实际国服开始时间确定，阈值为\(1706504400000\)毫秒。基础字典与建模活动分别处理，卡牌按国服发布时间过滤。数据manifest、partitions、excluded和unavailable文件保留质量排除及缺失，模型使用统一案例。多时间轴成员关系仅连接到分面明细，总体损失始终对物理案例去重。

T1500无训练及验证曲线，测试14期242前缀；T10的测试为13期248前缀。T1000测试16期281前缀，T2000为16期279前缀，原两线合计560；19线总体为5220。不同排名覆盖及每期前缀数造成活动内排名权重差别，\texttt{event\_rank\_balanced\_mape\_pct}另将活动内每条可见排名等权。该敏感性字段与正文活动内案例等权口径分别保存。

\section{固定方法、初始化与反馈敏感性}
\subsection{54种方法的完整成绩}
\begingroup\footnotesize\setlength{\tabcolsep}{3pt}
\begin{longtable}{p{6.8cm} r r r r}
\caption{完整54种固定规则及全部反馈策略}\label{tab:fullmenu}\\
\toprule
方法标识 & 验证MAPE & 测试MAPE & case-MAPE & P90 \\
\midrule\endfirsthead
\toprule
方法标识 & 验证MAPE & 测试MAPE & case-MAPE & P90 \\
\midrule\endhead
\path{uncorrected} & 6.672 & 7.082 & 6.929 & 17.567 \\
\path{ewma_reference} & 6.502 & 6.671 & 6.464 & 16.760 \\
\path{legacy_pair} & 6.556 & 6.540 & 6.342 & 16.559 \\
\path{time_proportion} & 18.985 & 17.699 & 17.306 & 41.982 \\
\path{history_growth} & 18.051 & 17.915 & 17.615 & 44.236 \\
\path{residual_mean} & 6.186 & 7.122 & 6.942 & 17.870 \\
\path{window5} & 6.464 & 6.697 & 6.533 & 16.859 \\
\path{window10} & 6.251 & 6.706 & 6.521 & 16.676 \\
\path{ewma_fast} & 6.888 & 7.137 & 6.915 & 18.076 \\
\path{pair_anchor_only} & 6.593 & 7.305 & 7.117 & 18.529 \\
\path{pair_difference_only} & 6.561 & 6.511 & 6.323 & 16.491 \\
\path{pair_swapped_anchor} & 6.566 & 6.926 & 6.711 & 17.113 \\
\path{same_filter_pair} & 6.502 & 6.671 & 6.464 & 16.760 \\
\path{same_panel_independent} & 6.502 & 6.671 & 6.464 & 16.760 \\
\path{kalman_level} & 6.591 & 6.762 & 6.560 & 16.875 \\
\path{kalman_trend} & 7.261 & 7.338 & 7.146 & 18.352 \\
\path{dlm_fullcov} & 6.510 & 6.672 & 6.462 & 16.750 \\
\path{pair_anchor_alpha_0p1} & 6.558 & 6.514 & 6.322 & 16.385 \\
\path{pair_anchor_alpha_0p3} & 6.556 & 6.561 & 6.360 & 16.566 \\
\path{pair_difference_alpha_0p3} & 6.503 & 6.623 & 6.420 & 16.738 \\
\path{pair_difference_alpha_0p9} & 6.603 & 6.535 & 6.336 & 16.563 \\
\path{pair_anchor_shrinkage_0} & 6.559 & 6.564 & 6.364 & 16.646 \\
\path{pair_anchor_shrinkage_5} & 6.557 & 6.551 & 6.352 & 16.613 \\
\path{pair_anchor_shrinkage_20} & 6.555 & 6.526 & 6.330 & 16.539 \\
\path{pair_difference_shrinkage_0} & 6.573 & 6.526 & 6.329 & 16.478 \\
\path{pair_difference_shrinkage_2p5} & 6.563 & 6.533 & 6.336 & 16.491 \\
\path{pair_difference_shrinkage_10} & 6.545 & 6.560 & 6.362 & 16.659 \\
\path{pair_minimum_pairs_3} & 6.556 & 6.540 & 6.342 & 16.559 \\
\path{pair_minimum_pairs_5} & 6.556 & 6.540 & 6.342 & 16.559 \\
\path{pair_minimum_pairs_10} & 6.556 & 6.540 & 6.342 & 16.559 \\
\path{pair_T400_T500} & 6.486 & 6.683 & 6.474 & 16.726 \\
\path{pair_T1000_T3000} & 6.555 & 6.581 & 6.380 & 16.615 \\
\path{pair_T2000_T3000} & 6.493 & 6.671 & 6.466 & 16.751 \\
\path{ewma_reference__delay24h} & 6.502 & 6.671 & 6.464 & 16.760 \\
\path{legacy_pair__delay24h} & 6.556 & 6.540 & 6.342 & 16.559 \\
\path{dlm_fullcov__delay24h} & 6.510 & 6.672 & 6.462 & 16.750 \\
\path{ewma_reference__delay7d} & 6.520 & 6.595 & 6.422 & 16.121 \\
\path{legacy_pair__delay7d} & 6.623 & 6.465 & 6.301 & 16.027 \\
\path{dlm_fullcov__delay7d} & 6.528 & 6.606 & 6.431 & 16.083 \\
\path{ewma_reference__delay_one_query} & 6.520 & 6.633 & 6.464 & 16.388 \\
\path{legacy_pair__delay_one_query} & 6.623 & 6.518 & 6.360 & 16.222 \\
\path{dlm_fullcov__delay_one_query} & 6.528 & 6.645 & 6.474 & 16.345 \\
\path{ewma_reference__drop10pct} & 6.547 & 6.676 & 6.483 & 16.760 \\
\path{legacy_pair__drop10pct} & 6.600 & 6.535 & 6.352 & 16.652 \\
\path{dlm_fullcov__drop10pct} & 6.564 & 6.671 & 6.474 & 16.798 \\
\path{ewma_reference__drop30pct} & 6.411 & 6.713 & 6.530 & 16.980 \\
\path{legacy_pair__drop30pct} & 6.418 & 6.555 & 6.382 & 16.749 \\
\path{dlm_fullcov__drop30pct} & 6.417 & 6.727 & 6.538 & 17.032 \\
\path{ewma_reference__missing_anchor2000} & 6.499 & 6.684 & 6.481 & 16.789 \\
\path{legacy_pair__missing_anchor2000} & 6.499 & 6.684 & 6.481 & 16.789 \\
\path{dlm_fullcov__missing_anchor2000} & 6.469 & 6.968 & 6.764 & 17.585 \\
\path{ewma_reference__missing_target1000} & 6.575 & 7.112 & 6.922 & 18.293 \\
\path{legacy_pair__missing_target1000} & 6.572 & 7.104 & 6.916 & 18.297 \\
\path{dlm_fullcov__missing_target1000} & 6.571 & 7.104 & 6.918 & 18.293 \\
\bottomrule
\end{longtable}
\noindent 全部误差单位为\%。延迟/删减只作用于残差校正反馈，基础预测沿用原V3缓存。每方法测试均为16期5220唯一案例；总体评分未连接时间条件成员关系。\par
\endgroup

\texttt{time\_proportion}使用当前分数除实际进度作为预测；\texttt{history\_growth}对同线合法历史，以进度高斯权重（带宽0.15）先计算每活动的对数剩余增长均值，再对活动等权平均；无同线历史时回退时间比例。两者均执行共同外部分数边界。它们提供简单增长控制，基础预测来源有别于残差校正方法。均值、5期或10期窗口与快EWMA分别改变记忆方式；其余基础树预测保持固定。

仅锚点消融对T1000和T2000都扣除\(A\)；仅差异消融对T1000扣除\(D\)，对T2000偏差设零。交换锚点将T1000作为参照并相应重写差异方向。两种同滤波负控共同采用完整配对面板、\(\alpha=.3\)、收缩常数10，最后才执行总偏差截断。同滤波分解与同面板独立输出的对应比较条件完全相同。

单因素扰动固定锚点\(\alpha\in\{.1,.3\}\)、差异\(\alpha\in\{.3,.9\}\)，锚点收缩\(k\in\{0,5,20\}\)、差异收缩\(k\in\{0,2.5,10\}\)和最低配对支持\(n_{\min}\in\{3,5,10\}\)，逐次只替换该项。替代线对固定为T400/T500、T1000/T3000和T2000/T3000。表中的最低支持扰动在当前留出期可能给出相同预测，其作用是记录实际触发情况。

\subsection{完整协方差DLM的训练期参数}
真实DLM在原两线坐标使用二维随机游走、零初始均值。各阶段以训练期已释放的完整配对外推误差估计参数：令\(C\)为误差水平样本协方差，\(V_\Delta\)为相邻误差差分的样本协方差，设
\begin{equation}
\begin{aligned}
R&=\operatorname{diag}\!\left[\max\{.5\operatorname{diag}(V_\Delta),10^{-6}\}\right],\\
Q&=\operatorname{PSD}_{10^{-6}}\{.1[.5C+.5\operatorname{diag}(C)]\},\\
P_0&=\operatorname{PSD}_{10^{-6}}(C+10^{-6}I).
\end{aligned}
\end{equation}
\(\operatorname{PSD}\)将对称矩阵的特征值下限设为\(10^{-6}\)。配对少于8期时，\(Q=.0001I,R=.01I,P_0=.01I\)。各反馈策略在首次验证前分别固定其训练期参数；后续只更新状态，不重估\(Q,R\)。训练期逐期诊断每次仍只使用该期合法训练历史。

滤波使用预测协方差加\(Q\)、创新和Kalman增益，并用Joseph形式更新协方差。缺排名时删去对应观测行，有效信息不足时保持预测步骤。坐标变换控制把均值、\(P_0,Q,R\)及观测载荷同时变换，测试核对两种坐标的结果等价。该传统方法的结构和参数估计选择均在新评分前固定。

合成DLM的固定参数另为\(Q_{a,d}=\operatorname{diag}(.0002,.002)\)、原始\(R_\theta=.01I\)、原始\(P_\theta=.1I\)，避免与真实训练估参混合。删除交叉项控制只删除变换后\(R\)的非对角项，其余设置相同。oracle使用真实生成\(Q,R\)及初始状态，拥有额外信息，排除于真实方法选择。

\subsection{反馈通道与到达时间假设}
基准V3用\(E_j<S_e\)并核对协调器已释放身份。额外24小时、7天和一次查询延迟仅删减校正器获得的残差行；基础树仍是原V3缓存。10\%和30\%删减由活动、排名、阶段与固定种子42的哈希确定，所有方法使用同一掩码；缺线情景分别删除T2000或T1000的残差反馈。

这些场景衡量给定基础预测的校正通道变化。24小时在本日程下没有改变三种主规则预测；7天和一次查询延迟可使部分损失下降，丢失也呈不同方向。缺少一整条线的反馈更直接检验回退与状态可识别性。历史实际采集时刻缺少完整日志，仿真延迟由设定时间规则生成，基础模型整体延迟重训属于后续独立协议。

\section{共享风险推导与补充诊断}
\subsection{部分共享、已有偏差及时间尺度}
按正文的部分共享平方风险推导，设\(D=d^2+v_1+v_2-2c>0\)。任务1的正向共享降低平方风险当且仅当\(v_1>c\)且\(0<w<2(v_1-c)/D\)，同时\(w\le1\)。受限最优权重为\(\operatorname{clip}\{(v_1-c)/D,0,1\}\)。若\(D=0\)，两估计的差在该目标下几乎处处为零，风险差为零。固定半共享\(w=1/2\)对任务1的条件是\(d^2<3v_1-v_2-2c\)，区别于两任务等权完全池化条件。

已有估计偏差分别为\(b_1,b_2\)时，混合偏差相对项为\([(1-w)b_1+w(b_2+d)]^2-b_1^2\)。随机任务差异与估计误差相关时，还须保留交叉项。这些条件使经验残差相关与风险参数的对应关系保持明确；高终分相关首先描述水平，同期残差相关又混合持久状态和瞬时冲击。

独立噪声方差为\(r\)的稳定EWMA方差为\(r\alpha/(2-\alpha)\)，其有效样本量\((2-\alpha)/\alpha\)在\(\alpha=.2,.3,.6\)时约为9、5.67、2.33。支持活动数\(n\)与有效窗口是两个量。随机游走过程方差\(q\)时，常增益一步状态预测MSE为
\begin{equation}
\frac{q+\alpha^2r}{\alpha(2-\alpha)},\qquad\frac qr=\frac{\alpha^2}{1-\alpha}\quad\text{（最优增益条件）}.
\end{equation}
因此0.2和0.6对应的过程/噪声比约为0.05和0.9，明确表达“锚点慢、差异快”的假设。实际配方还有支持收缩与边界，这些公式用于机制解释。

\subsection{共同支持、相关差及滞后筛查}
\begin{table}[!htbp]\centering
\caption{共同34期153线对的相关中位数差}\label{tab:corrcontrasts}
\small\setlength{\tabcolsep}{4pt}\renewcommand{\arraystretch}{1.15}
\begin{tabularx}{\linewidth}{L r C{4.3cm}}\toprule
比较 & 中位相关差 & 活动配对95\%区间 \\
\midrule
剩余增长−终分 & -0.357 & [-0.570, -0.166] \\
原始误差−终分 & -0.657 & [-0.802, -0.474] \\
校正误差−终分 & -0.642 & [-0.794, -0.457] \\
\bottomrule\end{tabularx}
\par\vspace{4pt}\begin{minipage}{\linewidth}\footnotesize 重采样对所有排名使用同一活动索引，避免把153个线对当独立样本。每次先重估相关再求中位数差；5000次、种子42。块3/5/7及基础路线分组均保留在原CSV，解释为探索性。\end{minipage}
\end{table}

共同面板153个线对的Spearman中位数依次为终分0.878、增长0.521、原始误差0.221、独立校正误差0.237。后两者的全局中位关系略有上升；T1000/T2000从原始0.283降至校正0.130，属于特定线对。正文按对象比较，保留这种异质性。

原外推误差历史含35期，滞后一阶仅34对；本轮34期共同面板另为控制半程增长支持而建立。既有循环位移与全局BY筛查\citep{benjamini2001by}未发现\(q<.05\)关系，离散置换水平、小样本与多重校正共同限制检验力。历史滤波还涉及多阶记忆与收缩，单个一阶相关覆盖其部分统计结构。

\subsection{合成网格、Monte Carlo单位及失败区域}
\upfigure{upgrade_synthetic_sharing}{fig:upgrade-synthetic}{首轮固定合成机制实验。16个生成条件、4种历史长度与9种方法完整保留；每格每长度300条完整Monte Carlo路径，跨格复用底层随机数。状态MSE相对独立EWMA的正差表示负迁移，oracle拥有额外信息。}{synthetic}
固定合成配置摘要为\path{45def3eee6388a994e8dbf22d79ec118f0e83d5c340b6475924a7da7bc9b56ba}。16格均从\(c=d=0\)开始，高斯创新及原始两线噪声跨期独立。各格复用标准正态随机数，不同方法共享路径，四种历史长度来自同一序列。独立Monte Carlo单位是每种子内的完整生成路径，每格每长度共300个；格子和历史长度之间的比较具有设计上的相关性。

\begin{table}[!htbp]\centering
\caption{预登记合成网格的完整方向计数}\label{tab:syntheticcount}
\small\setlength{\tabcolsep}{4pt}\renewcommand{\arraystretch}{1.15}
\begin{tabularx}{\linewidth}{L r r r}\toprule
固定方法 & 网格平均状态MSE & 状态风险下降/上升 & 模拟APE下降/上升 \\
\midrule
零校正 & 0.015699 & 16/48 & 16/48 \\
独立EWMA & 0.003454 & 参考 & 参考 \\
完全共享EWMA & 0.008202 & 32/32 & 32/32 \\
同滤波锚点负控 & 0.003454 & 64组数值等价 & 64组数值等价 \\
既有锚点差异 & 0.003652 & 19/45 & 21/43 \\
独立固定Kalman & 0.002791 & 47/17 & 46/18 \\
删除观测交叉项DLM & 0.003066 & 38/26 & 37/27 \\
完整协方差固定DLM & 0.002852 & 43/21 & 48/16 \\
已知生成参数oracle & 0.002285 & 64/0 & 64/0 \\
\bottomrule\end{tabularx}
\par\vspace{4pt}\begin{minipage}{\linewidth}\footnotesize 16格与4历史长度等权，64组合的权重由设计给定。状态MSE使用统一偏差截断；方向相对独立EWMA，负控按10的负12次方容差判数值等价。oracle拥有真实生成参数，其角色为诊断参照。\end{minipage}
\end{table}

低共同漂移、高差异漂移、低噪声且\(m=0\)时，历史16期的锚点差异状态风险减少0.002100，Monte Carlo区间[−0.002448,−0.001753]。相反，共同漂移0.0004、差异漂移0、噪声0.0064、\(m=0\)且历史32期时增加0.001561，区间[0.001235,0.001887]。载荷从\(m=0\)变为0.5时，锚点差异下降组合数从13/32变为6/32。所有负区域保留，参数未据其重调。

表格与原始记录分别保存截断后的\texttt{state\_mse}、截断前\texttt{raw\_state\_mse}和\texttt{synthetic\_unbounded\_ape\_pct}。模拟APE为\(100\operatorname{mean}|\exp(y_{\mathrm{next}}-b)-1|\)，只共享偏差截断，未加入真实游戏分数边界。全输出包含1728行总体、5184行逐种子及1728行配对结果，来源为\path{synthetic/summary.csv}、\path{per_seed.csv}与\path{paired_comparisons.csv}。

\section{新拟合学习曲线、种子与成本}
\subsection{首轮规模曲线与全部56单元}
每种规模按实际开始时间取最早16、32、48、67期，分别有576、1170、1766、2452个训练前缀，预测同一14期488个两线验证案例。规模、日历覆盖及活动组成共同变化。局部特征、增长目标和边界固定；CatBoost为200轮，共享拟合1个模型、独立拟合2个模型。神经训练60轮，三个种子最终分数均值为主。共享输出、专属头和完全独立的参数及优化预算随结构变化，表格保持原身份。
\begin{table}[!htbp]\centering
\caption{相同488验证案例上的训练规模诊断（MAPE / \%）}\label{tab:learning}
\small\setlength{\tabcolsep}{4pt}\renewcommand{\arraystretch}{1.15}
\begin{tabularx}{\linewidth}{r r r r r r}\toprule
训练活动 & 树共享 & 树独立 & 网络共享 & 网络专属头 & 网络独立 \\
\midrule
16 & 11.227 & 9.838 & 13.128 & 13.275 & 12.749 \\
32 & 8.765 & 8.093 & 12.010 & 11.155 & 10.196 \\
48 & 8.429 & 7.979 & 11.811 & 10.991 & 10.203 \\
67 & 6.912 & 6.312 & 10.949 & 10.198 & 9.314 \\
\bottomrule\end{tabularx}
\par\vspace{4pt}\begin{minipage}{\linewidth}\footnotesize 验证活动固定14期。网络三列为三种子最终分数平均后的评分；树固定种子42。当前局部特征和CatBoost200/神经基础拟合与19档回放的混合基础路线分别解释。\end{minipage}
\end{table}

\upfigure{upgrade_training_size}{fig:upgrade-learning}{首轮重新拟合的规模诊断。所有点反复使用同一14期488案例验证面板；训练规模与日历覆盖共同变化，网络参数与模型数随结构变化。二轮公平对照另在训练内部多终点执行。}{learning}
\begingroup\footnotesize\setlength{\tabcolsep}{3pt}
\begin{longtable}{r l l l r r}
\caption{新拟合学习曲线全部56个评分单元}\label{tab:learningseeds}\\
\toprule
训练期数 & 家族 & 结构 & 种子/集成 & MAPE / \% & P90 / \% \\
\midrule\endfirsthead
\toprule
训练期数 & 家族 & 结构 & 种子/集成 & MAPE / \% & P90 / \% \\
\midrule\endhead
16 & catboost & shared & 42 & 11.227 & 26.375 \\
16 & catboost & independent & 42 & 9.838 & 23.141 \\
16 & neural & shared & 42 & 13.099 & 29.416 \\
16 & neural & shared & 123 & 13.545 & 30.137 \\
16 & neural & shared & 2026 & 14.542 & 31.628 \\
16 & neural & shared & mean & 13.128 & 29.590 \\
16 & neural & rank\_heads & 42 & 14.239 & 31.371 \\
16 & neural & rank\_heads & 123 & 14.830 & 31.936 \\
16 & neural & rank\_heads & 2026 & 15.872 & 32.884 \\
16 & neural & rank\_heads & mean & 13.275 & 29.893 \\
16 & neural & independent & 42 & 12.486 & 28.181 \\
16 & neural & independent & 123 & 13.141 & 28.446 \\
16 & neural & independent & 2026 & 14.145 & 32.337 \\
16 & neural & independent & mean & 12.749 & 28.842 \\
32 & catboost & shared & 42 & 8.765 & 21.060 \\
32 & catboost & independent & 42 & 8.093 & 19.528 \\
32 & neural & shared & 42 & 12.208 & 27.145 \\
32 & neural & shared & 123 & 11.794 & 24.650 \\
32 & neural & shared & 2026 & 13.302 & 28.270 \\
32 & neural & shared & mean & 12.010 & 26.569 \\
32 & neural & rank\_heads & 42 & 12.438 & 27.471 \\
32 & neural & rank\_heads & 123 & 10.794 & 24.491 \\
32 & neural & rank\_heads & 2026 & 12.020 & 26.711 \\
32 & neural & rank\_heads & mean & 11.155 & 24.901 \\
32 & neural & independent & 42 & 9.912 & 23.043 \\
32 & neural & independent & 123 & 10.346 & 23.616 \\
32 & neural & independent & 2026 & 11.407 & 25.249 \\
32 & neural & independent & mean & 10.196 & 23.515 \\
48 & catboost & shared & 42 & 8.429 & 21.129 \\
48 & catboost & independent & 42 & 7.979 & 19.681 \\
48 & neural & shared & 42 & 11.731 & 26.955 \\
48 & neural & shared & 123 & 11.554 & 24.686 \\
48 & neural & shared & 2026 & 13.267 & 28.109 \\
48 & neural & shared & mean & 11.811 & 26.174 \\
48 & neural & rank\_heads & 42 & 12.229 & 26.992 \\
48 & neural & rank\_heads & 123 & 10.776 & 25.102 \\
48 & neural & rank\_heads & 2026 & 11.789 & 26.683 \\
48 & neural & rank\_heads & mean & 10.991 & 25.609 \\
48 & neural & independent & 42 & 9.998 & 23.902 \\
48 & neural & independent & 123 & 10.350 & 24.069 \\
48 & neural & independent & 2026 & 11.168 & 24.806 \\
48 & neural & independent & mean & 10.203 & 23.563 \\
67 & catboost & shared & 42 & 6.912 & 16.563 \\
67 & catboost & independent & 42 & 6.312 & 15.613 \\
67 & neural & shared & 42 & 11.004 & 24.449 \\
67 & neural & shared & 123 & 10.359 & 21.352 \\
67 & neural & shared & 2026 & 12.600 & 27.363 \\
67 & neural & shared & mean & 10.949 & 23.753 \\
67 & neural & rank\_heads & 42 & 11.945 & 24.398 \\
67 & neural & rank\_heads & 123 & 9.683 & 21.354 \\
67 & neural & rank\_heads & 2026 & 11.370 & 25.710 \\
67 & neural & rank\_heads & mean & 10.198 & 22.286 \\
67 & neural & independent & 42 & 8.739 & 20.133 \\
67 & neural & independent & 123 & 9.646 & 21.949 \\
67 & neural & independent & 2026 & 10.816 & 24.150 \\
67 & neural & independent & mean & 9.314 & 20.987 \\
\bottomrule
\end{longtable}
\noindent mean指三个种子的最终解码分数均值，区别于三个误差的平均。全部行共用14期488案例。\par
\endgroup

\begingroup\footnotesize\setlength{\tabcolsep}{3pt}
\begin{longtable}{r l l r r r r}
\caption{实际拟合数、参数量及成本}\label{tab:fitbudget}\\
\toprule
训练期数 & 家族 & 结构 & 拟合数 & 每模型参数 & CPU秒 & 墙钟秒 \\
\midrule\endfirsthead
\toprule
训练期数 & 家族 & 结构 & 拟合数 & 每模型参数 & CPU秒 & 墙钟秒 \\
\midrule\endhead
16 & catboost & shared & 1 & — & 0.528 & 0.592 \\
16 & catboost & independent & 2 & — & 0.987 & 0.978 \\
16 & neural & shared & 3 & 1345 & 1.038 & 1.378 \\
16 & neural & rank\_heads & 3 & 1490 & 0.095 & 0.098 \\
16 & neural & independent & 6 & 1345 & 0.092 & 0.092 \\
32 & catboost & shared & 1 & — & 0.554 & 0.548 \\
32 & catboost & independent & 2 & — & 1.178 & 1.280 \\
32 & neural & shared & 3 & 1345 & 0.100 & 0.100 \\
32 & neural & rank\_heads & 3 & 1490 & 0.130 & 0.131 \\
32 & neural & independent & 6 & 1345 & 0.127 & 0.127 \\
48 & catboost & shared & 1 & — & 0.595 & 0.592 \\
48 & catboost & independent & 2 & — & 1.114 & 1.105 \\
48 & neural & shared & 3 & 1345 & 0.110 & 0.110 \\
48 & neural & rank\_heads & 3 & 1490 & 0.146 & 0.146 \\
48 & neural & independent & 6 & 1345 & 0.134 & 0.135 \\
67 & catboost & shared & 1 & — & 0.599 & 0.594 \\
67 & catboost & independent & 2 & — & 1.132 & 1.119 \\
67 & neural & shared & 3 & 1345 & 0.126 & 0.127 \\
67 & neural & rank\_heads & 3 & 1490 & 0.181 & 0.181 \\
67 & neural & independent & 6 & 1345 & 0.160 & 0.160 \\
\bottomrule
\end{longtable}
\noindent CatBoost各200轮，神经各60轮。网络独立每个种子拟合两个模型，参数列为单模型参数，合计预算随模型数变化。首轮初始化开销计入实际时间；这些小规模CPU成本具有环境依赖。\par
\endgroup

完全独立网络为两条排名各自拟合编码器和输出参数；训练预处理按选定子集拟合，参数张量互不共享。独立回归检查修改T2000监督时，T1000输出保持完全一致。各拟合范围内部按活动总权重相同，独立模型使用本线监督。三种子集成对最终解码分数求平均，所有检查点重新加载后核对封存预测。

网络共享输出、专属头与完全独立结构的单模型参数及拟合数见表~\ref{tab:fitbudget}。相同60轮保留训练时长规则，模型数与参数量随结构变化。该规模诊断不主张总计算量相等；实际墙钟和CPU时间用于复核已执行预算。

\subsection{同验证活动的配对差与解释范围}
\begin{table}[!htbp]\centering
\caption{两线新拟合神经集成的活动配对比较}\label{tab:learningcontrasts}
\small\setlength{\tabcolsep}{4pt}\renewcommand{\arraystretch}{1.15}
\begin{tabularx}{\linewidth}{r L r C{3cm} C{3cm}}\toprule
训练活动 & 比较 & 差值 / 百分点 & 活动配对95\%区间 & 5活动块95\%区间 \\
\midrule
16 & 专属头−共享 & 0.147 & [-0.389, 0.628] & [-0.298, 0.747] \\
16 & 独立−共享 & -0.379 & [-1.519, 0.610] & [-1.170, 0.337] \\
16 & 独立−专属头 & -0.526 & [-1.440, 0.238] & [-1.118, -0.005] \\
32 & 专属头−共享 & -0.855 & [-1.626, -0.140] & [-1.390, -0.221] \\
32 & 独立−共享 & -1.814 & [-2.984, -0.724] & [-2.665, -0.876] \\
32 & 独立−专属头 & -0.959 & [-1.618, -0.361] & [-1.385, -0.541] \\
48 & 专属头−共享 & -0.820 & [-1.553, -0.090] & [-1.267, -0.322] \\
48 & 独立−共享 & -1.608 & [-2.712, -0.591] & [-2.350, -0.747] \\
48 & 独立−专属头 & -0.788 & [-1.434, -0.236] & [-1.193, -0.385] \\
67 & 专属头−共享 & -0.751 & [-1.653, 0.162] & [-1.401, -0.042] \\
67 & 独立−共享 & -1.635 & [-2.694, -0.606] & [-2.399, -0.748] \\
67 & 独立−专属头 & -0.884 & [-1.261, -0.491] & [-1.157, -0.604] \\
\bottomrule\end{tabularx}
\par\vspace{4pt}\begin{minipage}{\linewidth}\footnotesize 均以同一14期验证的活动损失重采样，5000次、种子42。该验证已被多轮使用，区间未校正搜索。参数数量和独立模型数不同，预算见表~\ref{tab:fitbudget}。\end{minipage}
\end{table}

67期训练时，独立网络减专属头的MAPE差为−0.884个百分点，活动配对区间[−1.261,−0.491]。较小训练规模时部分比较区间跨零。每个规模按实际开始时间取最早\(N\)期，增大规模同时延长训练覆盖时期并改变活动组成。学习曲线各点复用同一验证面板，训练子集嵌套，点间相关；参数量、模型数和优化差异仍需进一步匹配。此处的验证结果服务于结构诊断，新测试或生产选择另行登记。

\subsection{既有19线神经残差结构}
既有19线网络在同一基础预测与独立校正上学习残差增量。编码器两层宽度16，输出头宽度8和1；专属头按支持收缩，未见排名用共享输出。相关图使用已结束活动的参考误差，每对至少8期、每线至多3邻居，支持不足回退对数排名距离先验。AdamW学习率.002、权重衰减.05、丢弃率.1、60轮及种子42/123/2026固定，历史不足12期回退。

共享输出、专属头、静态相关线与阶段图注意力的历史测试MAPE为6.716\%、7.240\%、7.213\%、7.218\%。专属头相对共享输出差0.524个百分点，既有活动iid区间[0.261,0.844]。相关线增量差−0.027个百分点，区间跨零。其负结果限定于该轮样本、输入与预算，梯度冲突未被实测。
\upfigure{upgrade_neural_comparison}{fig:upgrade-neural}{既有神经残差结构的匹配比较。活动配对与5活动循环块区间沿用该轮独立统计，历史测试已被探索。参数规模与优化仍有变化，比较范围限于已记录结构。}{neural}

\subsection{历史静态基线与预训练模型}
\begin{table}[!htbp]\centering
\caption{既有V2两线静态实验的描述性汇总}\label{tab:historicalstatic}
\small\setlength{\tabcolsep}{4pt}\renewcommand{\arraystretch}{1.15}
\begin{tabularx}{\linewidth}{L r r r}\toprule
方法 & 训练 / \% & 验证 / \% & 测试 / \% \\
\midrule
排名共享CatBoost & 3.611 & 6.430 & 12.755 \\
检索残差模型 & 3.602 & 5.355 & 23.141 \\
双时钟增长网络 & 5.196 & 9.299 & 12.577 \\
进度增长网络 & 5.402 & 9.605 & 12.398 \\
后续模型组合 & 2.277 & 4.658 & 14.062 \\
TabPFNv2上下文增强 & 0.931 & 5.979 & 12.747 \\
TabPFNv2基础配方 & 2.246 & 6.598 & 13.074 \\
多时距进度辅助监督 & 5.587 & 9.262 & 12.129 \\
Chronos-2（零样本） & — & 33.842 & 31.131 \\
TimesFM2.5（零样本） & — & 38.057 & 35.125 \\
双线TCN三种子均值 & 3.000 & 5.560 & 14.392 \\
阶段路由组合 & 2.203 & 4.852 & 13.259 \\
\bottomrule\end{tabularx}
\par\vspace{4pt}\begin{minipage}{\linewidth}\footnotesize 训练列依方法包含训练拟合或上下文重代入；零样本方法留空。输入、权重版本和预算随历史配方变化。该表与V3回放具有不同信息条件。\end{minipage}
\end{table}

历史TabPFN的训练查询包含在训练上下文中，训练列属于上下文重代入；Chronos-2和TimesFM2.5为零样本推理，训练列留空。TimesFM的2024方法论文与实验2.5检查点分别记录，Chronos-2按2025预印本引用。各历史配方的输入表达、权重版本、原选中方法及指纹见原运行清单。跨协议或跨排名总分用于描述各自任务，同案例和同信息条件的增量另以匹配比较判断。

\section{区间目标、尾部与前瞻接口审计}
\subsection{经验区间与联合覆盖}
\begin{table}[!htbp]\centering
\caption{三种覆盖目标的区别}\label{tab:intervaltargets}
\small\setlength{\tabcolsep}{4pt}\renewcommand{\arraystretch}{1.15}
\begin{tabularx}{\linewidth}{r r r r r r}\toprule
名义 / \% & 组数 & 活动数 & 组内联合 / \% & 逐案例 / \% & 整活动联合 / \% \\
\midrule
80 & 878 & 16 & 79.86 & 93.01 & 0.00 \\
90 & 878 & 16 & 90.47 & 96.90 & 6.25 \\
\bottomrule\end{tabularx}
\par\vspace{4pt}\begin{minipage}{\linewidth}\footnotesize 校准目标是同活动、排名和阶段内的全部前缀。整活动联合要求该活动所有排名和阶段均覆盖，是更强的额外诊断；80\%和90\%各有0/16、1/16期整活动全部覆盖。各目标的块3/5/7区间另存完整CSV。\end{minipage}
\end{table}

每组分数为同活动、排名和阶段内最大的绝对中心对数误差。因此名义水平对应组内全部冻结前缀的联合覆盖。逐案例覆盖较高，覆盖对象较弱；整活动全部排名与阶段同时覆盖更强，当前构造未针对该目标校准。可用率、无界回退和有限支持宽度独立报告，避免用无界区间掩盖支持不足。

支持不足先回退到合法历史的活动阶段最大误差池。T1500可在没有自身历史时借用该池，其\texttt{own\_support}与实际采用的\texttt{support}各自保存。当前5220案例均为有限支持，80\%水平采用支持数最小4、90\%最小9。组内联合覆盖接近名义的回溯现象仍处于单一已探索面板；时间依赖及活动变化使交换性假设需要额外检验。块区间、逐排名/阶段/类型/乐队可靠性保留于完整统计表。

\section{训练35期的候选、消融与公平预算}\label{sec:round2supp}
\subsection{全部候选及支持守卫}
固定专家\(I\)为本线\(\alpha=.3,k=10\)独立EWMA；\(S\)在完整排名组内等权平均后采用相同滤波与收缩，内部偏差边界均为\(\pm.3\)。无完整组历史时\(S=I\)。同面板独立控制把\(I\)也限定到完整组活动，区分池化与支持改变。冻结两线面板中两种独立输出精确一致，原外部独立专家与本轮固定独立专家在11536案例的最大分数差也为零，两者来源身份各自保留。

风险系数只使用已发专家的已释放误差。每活动排名阶段内先对物理案例计算\(a_j=\operatorname{mean}h^2\)与\(c_j=\operatorname{mean}(\epsilon h)\)，再按活动权重汇总。全历史风险用\(k=4/12\)；近期风险用衰减.85、\(k=8\)。向零ridge收缩先于截断，避免改变估计量。逐活动删一后重新估计系数的最小至最大跨度仅描述敏感性。

门控最近最多12个完整活动，要求至少8对、原始误差Pearson\(\rho\ge.15\)、\(\mathrm{RMS}(y_1-y_2)/(\mathrm{RMS}(y_1)+\mathrm{RMS}(y_2)+10^{-8})\le.5\)，差分后的相应比例至多.75；常量序列相关按零处理。全部阈值固定。选择块30个活动阶段中仅1个开启，诊断块全部关闭。门控前风险系数及删除范围继续保存，避免将实际零系数与开关前估计混淆。
\begin{table}[!htbp]\centering
\caption{训练内滚动选择：全部12个固定共享候选}\label{tab:round2adaptive}
\small\setlength{\tabcolsep}{4pt}
\begin{tabularx}{\linewidth}{L r r r}\toprule
候选 & 选择10期 / \% & 末10期 / \% & 末段log平方风险 \\
\midrule
独立EWMA（预定选择） & 4.977 & 5.918 & 0.010006 \\
固定1/4共享 & 5.247 & 5.954 & 0.009912 \\
固定1/2共享 & 5.600 & 6.109 & 0.010134 \\
两线完全共享 & 6.534 & 6.816 & 0.011522 \\
风险收缩 k=4 & 5.028 & 5.968 & 0.010038 \\
风险收缩 k=12 & 5.015 & 5.956 & 0.010027 \\
近期风险 k=8 & 5.013 & 6.017 & 0.010095 \\
稳健门控 k=8 & 4.998 & 5.918 & 0.010006 \\
同面板独立 & 4.977 & 5.918 & 0.010006 \\
风险规则仅改T1000 & 5.008 & 5.918 & 0.010007 \\
风险规则仅改T2000 & 4.983 & 6.017 & 0.010094 \\
相邻排名风险规则 & 4.985 & 5.925 & 0.009974 \\
\bottomrule\end{tabularx}
\par\vspace{3pt}\begin{minipage}{\linewidth}\footnotesize 两线主面板；选择块372案例，末段352案例；各10活动等权。次10期按候选顺序处理并列，选中独立EWMA。log平方损失与MAPE分别保留。\end{minipage}
\end{table}

\begin{table}[!htbp]\centering
\caption{末10期的共享系数与删除敏感性}\label{tab:round2lambda}
\small\setlength{\tabcolsep}{4pt}
\begin{tabularx}{\linewidth}{L r r r r}\toprule
规则 & 正系数/状态 & 均值 & 最大值 & 最大删除跨度 \\
\midrule
风险收缩 k=4 & 15/60 & 0.033 & 0.398 & 0.625 \\
风险收缩 k=12 & 15/60 & 0.027 & 0.321 & 0.500 \\
近期风险 k=8 & 27/60 & 0.072 & 0.583 & 0.802 \\
稳健门控 k=8 & 0/60 & 0.000 & 0.000 & 0.802 \\
风险规则仅改T1000 & 6/60 & 0.009 & 0.152 & 0.365 \\
风险规则仅改T2000 & 21/60 & 0.063 & 0.583 & 0.802 \\
相邻排名风险规则 & 19/60 & 0.265 & 1.000 & 1.000 \\
\bottomrule\end{tabularx}
\par\vspace{3pt}\begin{minipage}{\linewidth}\footnotesize 60状态为10活动×2排名×3阶段。门控行的删除跨度属于门控前风险估计；实际系数全为零。删除跨度描述有限支持敏感性。\end{minipage}
\end{table}


相邻组按19档顺序固定成对，末尾单线保持独立；未用相关矩阵重排。T1000所在(1000,1500)组因T1500全缺始终回退，T2000的(2000,3000)组有支持，故该候选在主两线上具有非对称信息条件。B共35次固定专家计算、420次候选混合，无基础重训；候选、输出与选择seal完整保留。

\subsection{全部18项结构消融}
各规则共用相同基础预测、冻结案例与共同分数边界。单线消融把另一线留在独立参考；仅差异对T2000偏差设零，与仅替换T1000保留T2000独立的控制分别定义。统一\(\alpha\)、统一收缩、完整支持及逐项取消收缩按固定配置执行。
\begin{table}[!htbp]\centering
\caption{训练末10期：18项预定结构消融}\label{tab:round2ablations}
\small\setlength{\tabcolsep}{4pt}
\begin{tabularx}{\linewidth}{L r r}\toprule
固定规则 & 两线MAPE / \% & 相对独立差 / pp \\
\midrule
未校正 & 6.744 & 0.827 \\
独立EWMA & 5.918 & 0.000 \\
原锚点差异 & 6.388 & 0.471 \\
仅替换T1000 & 6.439 & 0.522 \\
仅替换T2000慢锚点 & 5.867 & -0.051 \\
T1000作锚，仅替换T2000 & 6.551 & 0.633 \\
仅锚点 & 7.035 & 1.117 \\
仅差异 & 6.355 & 0.437 \\
同滤波坐标负控 & 5.918 & 0.000 \\
同面板独立 & 5.918 & 0.000 \\
统一衰减系数 & 5.944 & 0.027 \\
统一收缩 & 6.321 & 0.403 \\
完整配对支持 & 6.388 & 0.471 \\
锚点不收缩 & 6.424 & 0.506 \\
差异不收缩 & 6.527 & 0.609 \\
全部不收缩 & 6.554 & 0.636 \\
锚点衰减改为.3 & 6.396 & 0.479 \\
差异衰减改为.3 & 5.923 & 0.006 \\
\bottomrule\end{tabularx}
\par\vspace{3pt}\begin{minipage}{\linewidth}\footnotesize 各规则共用10期352案例。方法标识对应封存配置；仅替换某一线时，另一线保留独立校正。该表保留预定消融身份，末段未用于重选。\end{minipage}
\end{table}

均值、跨期标准差、相邻期变化和lag1逐排名阶段保留于\path{C/runs/c1_002_scores/cross_event_stability.csv}及\path{event_rank_phase_residuals.csv}。稳定性量基于末10活动，阶段行保持同活动关联。原始误差、独立校正与原配方的变化共同报告，行为原因与统计机制分开解释。

\subsection{多终点历史选择、全部种子与实际预算}
N32full与N16span具有相同首末训练期，三个终点跨度分别297、305.875、341天；N16recent跨度145、153.875、175天。前缀数随冻结缺失变化。稀疏选择同时改变活动类型与组成；跨历史组编码的类别支持和\(d=56/57\)也变化。各单元自己的训练预处理对共享和独立结构完全相同，独立网络只接收所属排名监督。

共享网络无偏置\(d\to32\)编码器加两个带偏置输出头，总参数\(32d+66\)；两独立网络各为带偏置\(d\to16\to1\)，合计同为\(32d+66\)。参数相等按实际张量逐项核对。优化器为AdamW（学习率.002、权重衰减.05），全批次、无早停；活动等权且活动内排名等权。三种子集成按最终分数算术平均，非偏差或对数预测平均。

共享前向线性MAC为\(\mathrm{epochs}\,n(32d+64)\)，独立为\(\mathrm{epochs}\,n(16d+16)\)。共享每例计算两个头，独立每例只计算自己的头。相同60轮组的标量参数更新数相同，共享MAC约2.035倍；共享30轮的MAC比独立60轮多约1.73\%，每参数更新次数减半。CPU计时只覆盖拟合循环，初始化、输入读取、预处理和序列化成本另计。
\begin{table}[!htbp]\centering
\caption{公平训练的实际更新与计算预算}\label{tab:round2budget}
\small\setlength{\tabcolsep}{4pt}
\begin{tabularx}{\linewidth}{L r r r r}\toprule
结构 & 优化器步数 & 标量参数更新 & 前向线性MAC & 拟合CPU秒 \\
\midrule
independent60 & 3240 & 3044520 & 1175604480 & 0.563 \\
shared30 & 810 & 1522260 & 1195960320 & 0.249 \\
shared60 & 1620 & 3044520 & 2391920640 & 0.517 \\
\bottomrule\end{tabularx}
\par\vspace{3pt}\begin{minipage}{\linewidth}\footnotesize MAC只计训练前向线性层；反向传播、激活、偏置、优化器及预处理另有成本。共享30轮比独立60轮多约1.73\%前向MAC，每参数更新次数减半。\end{minipage}
\end{table}

\begingroup\footnotesize\setlength{\tabcolsep}{3pt}
\begin{longtable}{r p{2.5cm} r r r r}
\caption{训练期公平对照的全部27个种子单元}\label{tab:round2fairseeds}\\
\toprule
终点 & 历史组 & 种子 & 独立60轮 & 共享60轮 & 共享30轮 \\
\midrule\endfirsthead
\toprule
终点 & 历史组 & 种子 & 独立60轮 & 共享60轮 & 共享30轮 \\
\midrule\endhead
32 & N32full & 42 & 8.743 & 9.381 & 9.256 \\
32 & N32full & 123 & 7.868 & 8.804 & 8.135 \\
32 & N32full & 2026 & 8.711 & 8.138 & 7.850 \\
32 & N16span & 42 & 8.526 & 8.047 & 9.671 \\
32 & N16span & 123 & 8.001 & 9.307 & 9.803 \\
32 & N16span & 2026 & 7.933 & 9.309 & 9.522 \\
32 & N16recent & 42 & 8.986 & 8.943 & 9.545 \\
32 & N16recent & 123 & 7.502 & 8.814 & 7.595 \\
32 & N16recent & 2026 & 8.363 & 8.009 & 8.771 \\
44 & N32full & 42 & 8.705 & 7.135 & 8.324 \\
44 & N32full & 123 & 6.440 & 7.195 & 8.389 \\
44 & N32full & 2026 & 6.060 & 7.447 & 8.863 \\
44 & N16span & 42 & 10.731 & 10.080 & 9.386 \\
44 & N16span & 123 & 10.098 & 9.878 & 10.010 \\
44 & N16span & 2026 & 10.263 & 10.437 & 11.504 \\
44 & N16recent & 42 & 8.580 & 7.826 & 8.811 \\
44 & N16recent & 123 & 8.443 & 8.019 & 9.399 \\
44 & N16recent & 2026 & 6.694 & 7.073 & 8.590 \\
56 & N32full & 42 & 14.265 & 13.227 & 16.488 \\
56 & N32full & 123 & 11.271 & 11.295 & 14.290 \\
56 & N32full & 2026 & 10.697 & 11.340 & 16.488 \\
56 & N16span & 42 & 13.683 & 14.515 & 16.299 \\
56 & N16span & 123 & 11.920 & 11.828 & 15.111 \\
56 & N16span & 2026 & 11.891 & 13.073 & 16.494 \\
56 & N16recent & 42 & 13.959 & 10.867 & 12.726 \\
56 & N16recent & 123 & 9.662 & 10.589 & 12.674 \\
56 & N16recent & 2026 & 14.750 & 12.371 & 15.090 \\
\bottomrule
\end{longtable}
单位为活动等权MAPE（\%）。27个种子单元涉及108次物理拟合；共享结构各27次，独立结构54次。种子与重复终点维持活动重叠关系。\par
\endgroup


C2按结构、组、终点及种子保存MAE/RMSE/P90与配对活动差。12期汇总的总参数匹配差，独立减共享60轮在三历史组的iid区间均跨零。三个终点每个仅4期，较长块区间标注支持不足；合并12期的长块敏感性亦保留少量长段的精度限制。九单元及三个种子用于条件比较，活动数按12期计。

\section{受限机制参数、校准模拟与真实迁移}\label{sec:round2mechanism}
\subsection{原始矩、参数支持与适配缺口}
校准主面板为\(.25\le p<.9\)的T1000/T2000活动均值，前25期定参数、末10期诊断。差分矩、PSD投影、偏度峰度、JB筛查、方差前后比及滞后关联完整保存。JB为渐近筛查，前段108个检验中10个\(p<.05\)、2个BH\(q<.05\)，后段分别5和1个；样本9至25且有时间依赖，高斯适配的判断保留检验力范围。

中阶段原始\(Q\)含负特征值，\(R\)对角比为\DRatio，显示短序列波动与受限等噪声模型的适配缺口。完整900点网格固定\(m\in\{0,.5,1\}\)等离散值，描述性\(2\Delta\mathrm{NLL}\le10\)保留17点。最优共同过程项落下边界，噪声相关落上边界；支持点\(m=0\)的解释条件包括等观测方差及独立过程创新。真实共同或私有玩家状态的唯一识别仍需更多观测和更一般模型。
\begin{table}[!htbp]\centering
\caption{受限局部水平族的参数支持范围}\label{tab:round2calibration}
\small\setlength{\tabcolsep}{4pt}
\begin{tabularx}{\linewidth}{L r r r r r}\toprule
参数 & 最优 & 支持下限 & 支持上限 & 删一下限 & 删一上限 \\
\midrule
qc & 0.000000 & 0.000000 & 0.000736 & 0.000000 & 0.000000 \\
qd & 0.005888 & 0.002944 & 0.011776 & 0.002944 & 0.005888 \\
r & 0.000736 & 0.000147 & 0.001472 & 0.000736 & 0.000736 \\
rho & 0.500000 & -0.500000 & 0.500000 & 0.500000 & 0.500000 \\
m & 0.000000 & 0.000000 & 0.000000 & 0.000000 & 0.000000 \\
q\_r\_ratio & 4.000000 & 1.000000 & 45.000000 & 2.000000 & 4.000000 \\
\bottomrule\end{tabularx}
\par\vspace{3pt}\begin{minipage}{\linewidth}\footnotesize 900点离散网格；描述性支持条件为2ΔNLL≤10，保留17点。参数限于预定等观测方差与独立过程创新族。删一范围仅使用删除后所选参数；原表附带似然元数据属于完整25期。\end{minipage}
\end{table}


删一参数确由删除后似然选出，原\path{delete_one_profile.csv}附带的nll、support及样本量字段引用完整25期网格。\path{D_delete_one_metadata_audit.json}将其身份澄清为\texttt{original\_full25\_*}；本表删一范围仅使用删除后所选参数值。H另以标量滤波复算完整网格及一个固定删除点，参数选择一致。真实运行前的块接缝修订改为重抽原活动附着的真实差分贡献，未将拼接段的跳变当成真实相邻差分。

\subsection{24格条件模拟的完整失败区域}
低、中、高三个支持锚点的\(\operatorname{tr}Q/\operatorname{tr}R\)为1/4/45，三者均\(m=0\)；中位比锚点按事前顺序选实际支持点，身份与最优点分开。每锚点固定高斯、方差规范t5、5\%方差规范污染、共同突变、差异突变、1期延迟、3期延迟及异方差八种条件。每格三种子各100路径、长度35，末10期评分；共43200条路径方法记录。格间共享底层随机数，独立单位为每格的完整路径。
\begin{table}[!htbp]\centering
\caption{真实量纲校准模拟中的完整方向计数}\label{tab:round2sim}
\small\setlength{\tabcolsep}{4pt}
\begin{tabularx}{\linewidth}{L r r}\toprule
规则 & 状态MSE减少/增加 & 模拟APE减少/增加 \\
\midrule
完全共享 & 0/24 & 0/24 \\
固定半共享 & 0/24 & 0/24 \\
近期风险 & 9/15 & 9/15 \\
校准完整协方差DLM & 22/2 & 14/10 \\
已知基准协方差oracle & 22/2 & 22/2 \\
\bottomrule\end{tabularx}
\par\vspace{3pt}\begin{minipage}{\linewidth}\footnotesize 3个参数锚点×8扰动，共24格；每格3种子各100路径，方法与格间共用底层随机数。独立Monte Carlo单位为完整路径；24格描述设计覆盖。\end{minipage}
\end{table}


完全共享与半共享的24格配对Monte Carlo区间在状态MSE及模拟APE上均为正。近期风险在三个3期延迟格均增加状态风险，在共同突变格均减少相对独立风险。固定校准DLM的潜在状态MSE和模拟APE方向在部分格不同，真实得分边界也未进入模拟APE。oracle知道基准生成\(Q,R\)和初态，仍用固定局部水平递归，没有突变预知或逐时异方差信息；其扰动格负结果一并保留。

输出\texttt{raw\_state\_mse}指外部统一偏差截断之前，I/S专家内部\(\pm.3\)截断仍保留。\texttt{clipped\_fraction}只计额外外部截断；字段原名及原值未覆盖。二轮校准锚点均\(m=0\)，首轮手工网格中的载荷.5条件承担另外的机制范围。

\subsection{模拟后追加的固定DLM真实迁移}
追加计划在看过模拟后提出，再封存固定参数与代码；无额外参数搜索。仅对中阶段末10期120案例应用前25期最优\(Q,R,\mu\)及\(P_0=.1I\)，合法完整配对历史由25增至34，每期全部时点共用期初状态。H独立标量复算固定递归、当前分数边界及三比较规则，MAPE与原输出一致。

独立法为\DTransferI\%，固定DLM为\DTransfer\%，观测平方对数损失为\DTransferLog；其余逐期对照保留于\path{D/posthoc_transfer_run_001/event_comparisons.csv}。该运行逐期预测在内存收集后整体封存，未提供逐期真实公开发布时效证据。其结论覆盖固定信息集下的事后回溯迁移，正式前瞻结果由新的发布记录评价。

\section{条件可靠性、评分与未来样本量}\label{sec:round2reliability}
\subsection{完整组分层与宽度代价}
每名义水平5220案例、878完整组、16活动；同组所有冻结前缀同时覆盖才记组成功。\path{E/reliability_002/interval_conditions.csv}给76行逐案例条件，\path{interval_full_group_conditions.csv}给保持完整组的76行。排名、阶段、实际进度、类型、支持及区间来源全部保留，未据误差筛选活动。
\begin{table}[!htbp]\centering
\caption{预定条件审计的低覆盖与回退实例}\label{tab:round2conditions}
\small\setlength{\tabcolsep}{4pt}
\begin{tabularx}{\linewidth}{L r r r r}\toprule
条件 & 名义 / \% & 活动/组数 & 联合覆盖 / \% & 95\%活动区间 \\
\midrule
T1000 & 80 & 16/47 & 65.62 & [50.00,81.25] \\
T1000 & 90 & 16/47 & 87.50 & [72.92,97.92] \\
challenge类型 & 80 & 3/168 & 59.97 & [51.85,66.67] \\
challenge类型 & 90 & 3/168 & 74.79 & [62.96,84.21] \\
跨排名回退池 & 80 & 4/12 & 100.00 & [100.00,100.00] \\
跨排名回退池 & 90 & 10/27 & 100.00 & [100.00,100.00] \\
\bottomrule\end{tabularx}
\par\vspace{3pt}\begin{minipage}{\linewidth}\footnotesize 完整76个逐案例分层与76个完整组分层保留于原生CSV。条件间活动、排名、阶段及历史支持组成共同变化。回退池高覆盖伴随较宽区间。\end{minipage}
\end{table}


全活动联合0/16与1/16的固定损失bootstrap分别退化为[0,0]及活动重抽样分布；独立同分布活动二项参考另给Clopper--Pearson区间[0,20.5907\%]及[.1581\%,30.2321\%]。该二项参考另加活动独立假设。festival与medley各1期，bootstrap为单点；全成功低支持组的单点同样反映固定样本观察。支持与排名、时间和活动组成相关，条件差异解释保留这些混杂。

跨排名\texttt{all\_event\_phase\_pool}回退的完整组覆盖为100\%，80\%水平有4期12组，90\%有10期27组；其活动等权相对宽度约113\%/163\%，自身排名阶段来源约34\%/44\%。不同水平回退数量不同，较高覆盖与宽度代价共同报告。\texttt{own\_support}=0的11案例来自1期，各层小分母维持原值。

误覆盖率\(\alpha=1-\ell\)的标准Winkler评分\citep{gneiting2007proper}为
\begin{equation}
\operatorname{IS}_{\alpha}(l,u;Y)=(u-l)+\frac2\alpha(l-Y)\mathbf1(Y<l)+\frac2\alpha(Y-u)\mathbf1(Y>u).
\end{equation}
先在活动内平均再活动等权，80\%/90\%分别为\EWinklerEighty 万/\EWinklerNinety 万EP。按实际终分归一化的评分另作描述；标准分数尺度和名义水平的越界代价各自保留。最大组误差校准与逐案例区间评分回答不同目标，本轮没有按评分重新选宽度。

\subsection{规划模型、采集预算与推断条件}
功效方差只取训练35期已发legacy\_pair减独立EWMA活动MAPE差。主19线样本标准差约.079753个百分点，两线约.717920；两线规则只影响少数排名，19线平均降低差异方差。采用双侧零差检验、非中心t和正态iid卡方上95\%\(\sigma\)，再用\((1+\rho)/(1-\rho)\)作AR(1)长期方差膨胀敏感性。
\begin{table}[!htbp]\centering
\caption{训练差异驱动的两线样本量敏感性}\label{tab:round2power}
\small\setlength{\tabcolsep}{4pt}
\begin{tabularx}{\linewidth}{L r r r}\toprule
假设备择改善 / pp & ρ=0 & ρ=.3 & ρ=.6 \\
\midrule
0.25 & 151 & 279 & 597 \\
0.5 & 40 & 72 & 151 \\
1.0 & 12 & 20 & 40 \\
\bottomrule\end{tabularx}
\par\vspace{3pt}\begin{minipage}{\linewidth}\footnotesize 双侧α=.05、功效.9、训练35期正态iid卡方区间的上限σ；非中心t基于独立正态活动差，AR(1)长期方差膨胀仅为规划近似。该表规划零差检验。\end{minipage}
\end{table}


在改善.5个百分点、功效.9及\(\rho=.3\)时，主19线模型输出\ConfirmGlobalN\ 期，两线次面板为\ConfirmPairN\ 期；20期pilot在两线该条件下的模型功效约38.09\%。这些输出基于明确规划假设。60期工程采集下限与两线n72共同驱动按10向上取整的80期计划，主19线指标及双侧\(\alpha=.05\)保持不变。超过实用改善.5个百分点的移位零假设另在\path{power_usefulness_margin_grid.csv}列出，与零差检验分别规划。

未来主HAC区间采用时间排序活动差\(d_t\)：\(\widehat\gamma_h=n^{-1}\sum_{t=h+1}^n(d_t-\bar d)(d_{t-h}-\bar d)\)，\(\widehat V=\widehat\gamma_0+2\sum_{h=1}^L(1-h/(L+1))\widehat\gamma_h\)，并报告\(\bar d\pm z_{.975}\sqrt{\widehat V/n}\)。带宽\(L=\min(n-1,\max(1,\lceil n^{1/3}\rceil))\)固定。Bartlett协方差估计与弱依赖渐近条件来自既有研究\citep{newey1987hac}；零方差轨迹标记无显著性决策。iid配对t及块长1/3/5/7的固定损失敏感性另列，未来有限样本精度仍由真实依赖决定。

原20期pilot登记完整保留。新的\path{confirmation_plan.json}加独立\path{confirmation_plan_addendum.json}共同规定80期未激活队列；单独读取原JSON中的模型n4会遗漏采集补充。主比较、主指标、差值定义及双侧检验、完整活动人口、失败回退、缺失留存和固定报告时点均在激活前封存。当前未来活动为0，未来MAPE为空。


\section{版本绑定、失败记录与独立重建}
输入、源码、配置、依赖与预测分别封存。H独立用标量算术重建B全部22680个候选活动排名阶段状态，\(\lambda\)最大差为\(3.22\times10^{-15}\)；C实际张量参数和5670次优化器更新逐项核对。H的独立临时项目完整重放B35期，并重拟C2一个预定单元，与原预测逐例一致。原C源码快照的导入闭包缺口在可搬运包中补全；全新依赖安装环境仍保留自己的验收范围。

原B backend的物理案例别名守卫由独立适配器补充，原源码和封存预测保留；新入口为\path{integration/replay_guarded.py}，配合\path{sharing_input_guard.py}。H在冻结35期验证新守卫与原输出零差。D的原块接缝修订在真实运行前完成；删一参数附带的full25似然元数据用独立审计文件澄清，原表保留。E的statistics001为中断，reliability001因非中心t极尾NaN另存修订；本文数字统一来自statistics002/reliability002，失败目录随档案保留。

稿件内的一键数值重建显式传入仓库根，避免目录层级推断：
\begingroup\small
\begin{verbatim}
.venv/bin/python \
  <manuscript>/scripts/bind_round2_evidence.py \
  --repo-root <absolute-repository> \
  --paper-root <absolute-manuscript>
\end{verbatim}
\endgroup
该命令先重建原15张表与宏，再绑定本轮11张表与宏，逐项记录数值来源和SHA。它只读已封存分数，不拟合、重新评分或改变选择。图表由封存数值绘制，随后在稿件目录运行\texttt{make all}，使用XeLaTeX、BibTeX及交叉引用刷新生成完整稿与独立补充材料。\texttt{make verify}核对引用、字体及LaTeX警告；逐页渲染和接触表另存QA目录。

项目全套unittest已由PI按单线程串行执行，469项通过。二轮B15、C19、D19、E12及独立守卫5项夹具另计，共70项通过。E相对路径夹具首次从工作目录误调用的失败日志保留，按仓库根契约重跑通过，代码保持原版本。软件测试检验执行边界与实现，预测精度来自独立评分。

\subsection{尚待完成的外部证据}
当前真实未来活动为0，80期新队列等待新的激活封存。完整人口清单、准时预测、原始来源及实际释放日志仍依赖上游采集和独立审计。跨服务器和跨游戏验证、历史元数据时点快照与玩家层行为资料，是解释和推广当前结论所需的新增证据。所有历史分数保留测试探索、验证复用与回溯元数据限制；本轮35期研究亦保留回溯身份。

\ifdefined\PaperMainDocument\else
\FloatBarrier
\begingroup\raggedright\small
\bibliographystyle{elsarticle-num}
\bibliography{references}
\endgroup
\fi
\EndSupplementDocument
